%%% ============================================================
%%% Joint Closure of Two Binary Operations on Z/10Z:
%%% Chain Structure and a Normalizer Identity
%%%
%%% B.R. Sanders, M. Gish (2026)
%%% Target venue: Algebraic Combinatorics
%%% MSC: 20N02, 20N05
%%%
%%% Verification: 4core_verification.py archived at
%%% DOI 10.5281/zenodo.18852047 (verification scripts archive)
%%%
%%% This is the SEED-NARROW version. The bundled version (with
%%% closed-form fixed point, Galois quartic, and mixing-weight
%%% observations) is preserved in master/four_core_FINAL_bundled.tex
%%% and is the source for follow-on Papers 2 and 3.
%%% ============================================================

\documentclass[11pt,reqno]{amsart}

\usepackage{amsmath,amssymb,amsthm,mathtools}
\usepackage[margin=1in]{geometry}
\usepackage[colorlinks=true,linkcolor=blue,citecolor=blue,urlcolor=blue]{hyperref}
\usepackage{microtype}
\usepackage{array}

%% -------- theorem environments --------
\theoremstyle{plain}
\newtheorem{theorem}{Theorem}
\newtheorem{lemma}[theorem]{Lemma}
\newtheorem{proposition}[theorem]{Proposition}
\newtheorem{corollary}[theorem]{Corollary}

\theoremstyle{definition}
\newtheorem{definition}[theorem]{Definition}

\theoremstyle{remark}
\newtheorem*{remark}{Remark}
\newtheorem*{notation}{Notation}

%% -------- macros --------
\newcommand{\Z}{\mathbb{Z}}
\newcommand{\Q}{\mathbb{Q}}
\newcommand{\R}{\mathbb{R}}
\newcommand{\Zten}{\Z/10\Z}
\newcommand{\Cfour}{\mathcal{C}}
\newcommand{\fuseT}{\widehat{T}}
\newcommand{\fuseB}{\widehat{B}}
\newcommand{\CL}{\mathrm{CL}}
\newcommand{\F}{\mathbb{F}}

\title[Joint Closure on $\Zten$]{Joint Closure of Two Binary Operations on $\Zten$:
Chain Structure and a Normalizer Identity}

\author{B.~R.~Sanders}
\address{7Site LLC, Hot Springs, AR, USA}
\email{brayden@7site.io}

\author{M.~Gish}
\address{Independent Researcher}

\subjclass[2020]{Primary 20N02; Secondary 20N05}

\keywords{commutative magma, sub-magma chain, joint closure,
fuse-normalizer, polynomial identity}

\date{\today}

\begin{document}

\begin{abstract}
We study a pair of commutative binary operations
$T,B:\Zten\times\Zten\to\Zten$, given explicitly by their Cayley
tables. The operation $B$ is the $N=10$ instance of the family
$\CL_{N}$ of binary composition tables on $\Z/N\Z$ studied in the
companion paper \cite{SandersGish2026Sigma}; the operation $T$ is
a coarse projection of the same underlying construction
(separately specified by its Cayley table), with image
$\{0,3,4,7,8,9\}\subset\Zten$. The two tables disagree on $93$
of the $100$ cells of the full Cayley table (and on $12$ of the
$16$ cells of the four-element restriction introduced below). Despite this near-total pointwise disagreement,
the pair $(T,B)$ exhibits two exact structural identities. (i) The
sub-magmas of $\Zten$ that are jointly closed under both $T$ and
$B$ form a strict $8$-element chain in the inclusion order, with
sizes $\{1,4,5,6,7,8,9,10\}$; sizes $2$ and $3$ are forbidden as
joint closures. (ii) On the minimal non-trivial element of this
chain, the four-element subset $\Cfour=\{0,7,8,9\}$, the two
operations have a common fuse-normalizer
$Z_{T}(p)=Z_{B}(p)=(p_{0}+p_{7}+p_{8}+p_{9})^{2}$ for every
$p\in\R^{10}$ supported on $\Cfour$. The chain enumeration walks
the forward orbit of a permutation $\sigma$ on $\Zten$ with one
$\sigma$-fixed bridge step. The normalizer identity is a
polynomial identity in four variables, exact and independently
verifiable by symbolic expansion. Both identities are checked by a
deterministic verification script. This paper establishes
$(T,B)$-joint closure as a finite structural object; the
closed-form fixed-point dynamics of the convex-combination
iteration $\alpha\fuseT+(1-\alpha)\fuseB$, the Galois extension
governing the runtime ratio $r/\beta$, and the mixing-weight
observations at five sample $\alpha$-values are the subjects of
companion papers in preparation.
\end{abstract}

\maketitle

%% -------- section 1: introduction --------
\section{Introduction}\label{sec:intro}

This paper studies a small object: the family of sub-magmas of
$\Zten$ that are simultaneously closed under two specific binary
operations $T$ and $B$. We establish two exact structural
identities and pose them as a finite combinatorial-algebraic
anchor.

The two operations are presented as Cayley tables in
\S\ref{sec:setup}. They are not independent: $B$ is the $N=10$
instance of the family $\CL_{N}$ of binary composition tables on
$\Z/N\Z$ studied in the companion paper
\cite{SandersGish2026Sigma}, which establishes a sharp two-sided
non-associativity decay rate $N\sigma(N)\to 2$ for $\CL_{N}$ on
squarefree moduli. The operation $T$ is a coarse projection of
the same underlying construction, separately specified by its
Cayley table. Both tables fix the value $0$ at the boundary
($T(0,j)$ and $B(0,j)$ are determined by row-zero conventions
exhibited in Definition~\ref{def:tables}) and assign a
distinguished role to the value $7$: in $T$, the value $7$ acts
as a row absorber ($T(7,j)=7$ for all $j$); in $B$, the value
$7$ does not absorb pointwise but plays a structural role visible
through the joint-closure analysis of \S\ref{sec:chain}. The
present paper takes the pair $(T,B)$ at $N=10$ as given and asks:
what structural identities does the pairing produce?

The two tables disagree at $93$ of the $100$ cells of the full
Cayley table; on the four-element restriction $\Cfour=\{0,7,8,9\}$
(introduced below), they agree on exactly $4$ of the $16$ cells.
The non-associativity densities
$\sigma(M)=|\{(a,b,c)\in\Zten^{3}:M(a,M(b,c))\neq M(M(a,b),c)\}|/N^{3}$
of the two operations differ by a factor exceeding $3.7$:
$\sigma(T)\approx 0.128$ while $\sigma(B)\approx 0.498$. Despite
this surface dissimilarity, two exact structural identities link
them tightly.

\paragraph{The chain (Theorem~\ref{thm:chain}).} The sub-magmas
of $\Zten$ jointly closed under both $T$ and $B$, partially
ordered by inclusion, form a totally ordered $8$-element chain
\[
\{0\}\subset\Cfour\subset\Cfour\cup\{6\}\subset\Cfour\cup\{5,6\}\subset
\Cfour\cup\{4,5,6\}\subset\Cfour\cup\{3,4,5,6\}\subset
\Cfour\cup\{2,3,4,5,6\}\subset\Zten,
\]
where $\Cfour=\{0,7,8,9\}$. The chain is strict: no size-$2$ or
size-$3$ subset of $\Zten$ is jointly closed. Among all $1023$
non-empty subsets of $\Zten$, exactly these eight satisfy joint
closure. The chain enumeration has a clean structural reading
(Remark~\ref{rem:sigma-walk}): after the size-$1\to 4$ jump
$\{0\}\to\Cfour$, each subsequent size adds one element, and the
sequence of added elements walks the forward orbit of a
permutation $\sigma$ on $\Zten$ with one $\sigma$-fixed bridge
step. Verification of the chain enumeration over $\F_p$ for
several primes $p$, indicating a $p$-universal persistence of the
structure, is a target of future work.

\paragraph{The 4-core normalizer identity (Theorem~\ref{thm:norm}).}
On the minimal non-trivial element $\Cfour$ of the chain, the
fuse-normalizer
\[
Z_{M}(p)\;=\;\sum_{c\in\Zten}\,\sum_{(i,j)\in\Zten^{2}\,:\,M(i,j)=c}p_{i}p_{j}
\]
is identical for $M=T$ and $M=B$:
\[
Z_{T}(p)\;=\;Z_{B}(p)\;=\;(p_{0}+p_{7}+p_{8}+p_{9})^{2}\qquad\text{for every }p\text{ supported on }\Cfour,
\]
even though $T$ and $B$ disagree on $12$ of the $16$ cells of the
$4\times 4$ restriction. The identity is a polynomial identity in
the four variables $(p_{0},p_{7},p_{8},p_{9})$ and arises by a
global cancellation across cells, not by pointwise equality.

\paragraph{Positioning.} The non-associativity rate
$N\sigma(N)\to 2$ for the family of binary tables containing $B$
is the subject of the companion paper
\cite{SandersGish2026Sigma}. The classical Drápal--Kepka tradition
of low-density associativity in quasigroups
\cite{Kepka1980,DrapalKepka1985} treats non-associativity counts
on cancellative structures; the present paper studies the joint
closure of two non-cancellative magmas (one of which sits in the
family of \cite{SandersGish2026Sigma}, the other its threshold
projection from a common substrate) and is not implied by the
density results in those references.

\paragraph{Forward directions.} The pair $(T,B)$ exhibits
substantial additional structure beyond the chain and normalizer
identity, including a closed-form algebraic fixed point of the
convex-combination iteration $\alpha\fuseT+(1-\alpha)\fuseB$ at
$\alpha=\tfrac12$, a quartic Galois extension governing the
runtime ratios, and mixing-weight observations at sample
$\alpha$-values. These results are presented in companion papers
in preparation \cite{SandersGish2026FixedPoint,SandersGish2026MixingWeight};
this paper establishes the chain and normalizer identity as the
finite structural anchor on which those results rest. A summary
of the forward agenda appears in \S\ref{sec:forward}.

\paragraph{Outline.} \S\ref{sec:setup} fixes the two tables and
defines the fuse and normalizer. \S\ref{sec:chain} proves the
joint-closure chain (Theorem~\ref{thm:chain}). \S\ref{sec:4core}
displays the size-$4$ shell $\Cfour$ as a corollary.
\S\ref{sec:norm} proves the normalizer coincidence
(Theorem~\ref{thm:norm}). \S\ref{sec:forward} sketches the
forward agenda. \S\ref{sec:verify} catalogs the verification
script, and \S\ref{sec:scope} states the limitations of the
present results.

%% -------- section 2: setup --------
\section{Definitions and the two tables}\label{sec:setup}

We work throughout on the additive group $\Zten$, identified with
the labels $\{0,1,\dots,9\}$.

\begin{definition}[The two tables]\label{def:tables}
The two binary operations $T,B:\Zten\times\Zten\to\Zten$ are
given by their Cayley tables:
\[
T = \left[
\begin{array}{*{10}{c}}
0 & 0 & 0 & 0 & 0 & 0 & 0 & 7 & 0 & 0 \\
0 & 7 & 3 & 7 & 7 & 7 & 7 & 7 & 7 & 7 \\
0 & 3 & 7 & 7 & 4 & 7 & 7 & 7 & 7 & 9 \\
0 & 7 & 7 & 7 & 7 & 7 & 7 & 7 & 7 & 3 \\
0 & 7 & 4 & 7 & 7 & 7 & 7 & 7 & 8 & 7 \\
0 & 7 & 7 & 7 & 7 & 7 & 7 & 7 & 7 & 7 \\
0 & 7 & 7 & 7 & 7 & 7 & 7 & 7 & 7 & 7 \\
7 & 7 & 7 & 7 & 7 & 7 & 7 & 7 & 7 & 7 \\
0 & 7 & 7 & 7 & 8 & 7 & 7 & 7 & 7 & 7 \\
0 & 7 & 9 & 3 & 7 & 7 & 7 & 7 & 7 & 7
\end{array}
\right],\qquad
B = \left[
\begin{array}{*{10}{c}}
0 & 1 & 2 & 3 & 4 & 5 & 6 & 7 & 8 & 9 \\
1 & 2 & 3 & 4 & 5 & 6 & 7 & 2 & 6 & 6 \\
2 & 3 & 3 & 4 & 5 & 6 & 7 & 3 & 6 & 6 \\
3 & 4 & 4 & 4 & 5 & 6 & 7 & 4 & 6 & 6 \\
4 & 5 & 5 & 5 & 5 & 6 & 7 & 5 & 7 & 7 \\
5 & 6 & 6 & 6 & 6 & 6 & 7 & 6 & 7 & 7 \\
6 & 7 & 7 & 7 & 7 & 7 & 7 & 7 & 7 & 7 \\
7 & 2 & 3 & 4 & 5 & 6 & 7 & 8 & 9 & 0 \\
8 & 6 & 6 & 6 & 7 & 7 & 7 & 9 & 7 & 8 \\
9 & 6 & 6 & 6 & 7 & 7 & 7 & 0 & 8 & 0
\end{array}
\right].
\]
Both $T$ and $B$ are commutative.
\end{definition}

\begin{notation}
We refer to the indices $0,7,8,9$ frequently; they form the
four-element shell $\Cfour=\{0,7,8,9\}$ that is the minimal
non-trivial jointly-closed sub-magma (Theorem~\ref{thm:chain}). In
proofs over $\Cfour$ we sometimes write
$(p_{0},p_{7},p_{8},p_{9})=(v,h,\beta,r)$.
\end{notation}

\begin{remark}[Origin of the pair]
The pair $(T,B)$ is not independent. The operation $B$ is the
$N=10$ instance of the family $\CL_{N}$ of binary composition
tables on $\Z/N\Z$ studied in the companion paper
\cite{SandersGish2026Sigma}, which establishes the
non-associativity decay rate $N\sigma(N)\to 2$ for that family on
squarefree moduli. The operation $T$ is a coarse projection of
the same underlying construction, separately specified by the
Cayley table of Definition~\ref{def:tables}. Concretely, $T$ has
two designated absorbing values: the boundary value $0$, on which
$T$ is mostly absorbing ($T(0,j)=0$ for $j\neq 7$), and the
absorber $7$, on which $T$ is fully absorbing ($T(7,j)=7$ for all
$j$). Outside the absorbing rows and columns, $T$ takes the
absorber value $7$ at most cells, with $10$ cells (the
\emph{residue cells}, forming five symmetric pairs) preserving
non-absorbing values in $\{3,4,8,9\}\subset\Zten$. The image of
$T$ is therefore $\{0,3,4,7,8,9\}$, a six-element subset of
$\Zten$, while the image of $B$ is the full $\Zten$. The present
paper takes the explicit pair $(T,B)$ at $N=10$ as given and
asks: what structural identities does the pairing produce?
\end{remark}

\begin{definition}[Sub-magma; joint closure]\label{def:submagma}
A subset $S\subseteq\Zten$ is a \emph{sub-magma} of
$(\Zten,M)$ if $M(i,j)\in S$ for all $i,j\in S$. We call $S$
\emph{jointly closed} (under $T$ and $B$) if it is a sub-magma
of both. The set of jointly closed subsets, ordered by inclusion,
is denoted $\mathcal{L}(T,B)$.
\end{definition}

\begin{definition}[The fuse and normalizer]\label{def:fuse}
For a binary operation $M:\Zten\times\Zten\to\Zten$ and a vector
$p=(p_{0},\dots,p_{9})\in\R^{10}$, the \emph{fuse}
$\widehat{M}(p)\in\R^{10}$ is the convolution along the
operation:
\[
\widehat{M}(p)_{c}\;=\;\sum_{(i,j)\in\Zten^{2}\,:\,M(i,j)=c}p_{i}p_{j},\qquad c\in\Zten.
\]
The \emph{normalizer} is
$Z_{M}(p)=\sum_{c\in\Zten}\widehat{M}(p)_{c}$.
\end{definition}

\begin{remark}
For any $p\in\R^{10}$, $Z_{M}(p)=\bigl(\sum_{i}p_{i}\bigr)^{2}$
because the fuse is bilinear and each ordered pair $(i,j)$
contributes $p_{i}p_{j}$ to exactly one output bucket. The
substantive content of Theorem~\ref{thm:norm} below is that for
$p$ supported on the four-element subset $\Cfour$, the per-bucket
sums $\sum_{c\in\Cfour}\widehat{M}(p)_{c}$ for the two operations
$M=T$ and $M=B$ are not merely equal in total mass but are
\emph{individually} equal as polynomials in $(p_{0},p_{7},p_{8},p_{9})$
to $(p_{0}+p_{7}+p_{8}+p_{9})^{2}$.
\end{remark}

%% -------- section 3: chain --------
\section{The joint-closure chain}\label{sec:chain}

\begin{theorem}[Joint-closure chain]\label{thm:chain}
The sub-magmas of $\Zten$ jointly closed under both $T$ and $B$
form a strict chain of length $8$ in inclusion order:
\[
\mathcal{L}(T,B)\;=\;\bigl\{S_{1}\subset S_{4}\subset S_{5}\subset
S_{6}\subset S_{7}\subset S_{8}\subset S_{9}\subset S_{10}\bigr\},
\]
where
\begin{align*}
S_{1}&=\{0\}, & S_{4}&=\{0,7,8,9\}, \\
S_{5}&=\{0,6,7,8,9\}, & S_{6}&=\{0,5,6,7,8,9\}, \\
S_{7}&=\{0,4,5,6,7,8,9\}, & S_{8}&=\{0,3,4,5,6,7,8,9\}, \\
S_{9}&=\{0,2,3,4,5,6,7,8,9\}, & S_{10}&=\Zten.
\end{align*}
The chain is strict: no subset of $\Zten$ of size $2$ or size $3$
is jointly closed.
\end{theorem}

\begin{proof}
We verify the joint-closure status of each candidate subset by
direct computation. The proof has two parts: (a) confirming that
each $S_{k}$ in the displayed list is jointly closed, and (b)
confirming that no other non-empty subset is jointly closed.

\textbf{Part (a) — the eight subsets are jointly closed.} For
each $k\in\{1,4,5,6,7,8,9,10\}$, the verification reduces to
checking that $T(i,j)\in S_{k}$ and $B(i,j)\in S_{k}$ for all
$i,j\in S_{k}$. We display the size-$4$ case explicitly and refer
to the verification script (\S\ref{sec:verify}) for the others.

For $S_{4}=\{0,7,8,9\}$, the $4\times 4$ restrictions are
\[
T\big|_{S_{4}\times S_{4}}=\begin{pmatrix}
0 & 7 & 0 & 0 \\
7 & 7 & 7 & 7 \\
0 & 7 & 7 & 7 \\
0 & 7 & 7 & 7
\end{pmatrix},
\qquad
B\big|_{S_{4}\times S_{4}}=\begin{pmatrix}
0 & 7 & 8 & 9 \\
7 & 8 & 9 & 0 \\
8 & 9 & 7 & 8 \\
9 & 0 & 8 & 0
\end{pmatrix}.
\]
Every entry of $T\big|_{S_{4}\times S_{4}}$ lies in
$\{0,7\}\subseteq S_{4}$. Every entry of
$B\big|_{S_{4}\times S_{4}}$ lies in $\{0,7,8,9\}=S_{4}$.

\textbf{Part (b) — no other subset is jointly closed.} We
dispatch sizes $2$ and $3$ by exhaustive case analysis on the
diagonal of $B$, and sizes $4$ through $9$ by enumeration over
the relevant size class. Inspecting the diagonal:
\[
B(j,j)=\begin{cases}
0 & j=0,\\
j+1 & j\in\{1,2,3,4,5,6,7\},\\
7 & j=8,\\
0 & j=9.
\end{cases}
\]

\emph{Size $2$.} Let $S=\{a,b\}$ with $a<b$. Joint closure
requires $B(a,a),B(b,b)\in S$ and $T(a,a),T(b,b)\in S$, plus the
off-diagonal constraints.

\emph{Subcase $0\in S$.} Then $S=\{0,a\}$ with $a\in\{1,\dots,9\}$.
The diagonal $B(a,a)$ must lie in $\{0,a\}$. From the diagonal
table: for $a\in\{1,\dots,7\}$, $B(a,a)=a+1\notin\{0,a\}$; for
$a=8$, $B(8,8)=7\notin\{0,8\}$; for $a=9$, $B(9,9)=0\in\{0,9\}$
and $B(0,9)=B(9,0)=9\in\{0,9\}$, so $\{0,9\}$ is $B$-closed.
However $T(9,9)=7\notin\{0,9\}$, so $\{0,9\}$ fails $T$-closure.

\emph{Subcase $0\notin S$.} Then $S\subseteq\{1,\dots,9\}$, and
both $B(a,a)\in S$ and $B(b,b)\in S$ are required. Since
$B(j,j)\neq j$ for all $j\in\{1,\dots,9\}$, we need $B(a,a)=b$
and $B(b,b)=a$. The only unordered pair
$\{a,b\}\subseteq\{1,\dots,9\}$ satisfying this is $\{7,8\}$
(since $B(7,7)=8$ and $B(8,8)=7$). However $B(7,8)=9\notin\{7,8\}$,
so $\{7,8\}$ fails $B$-closure. The remaining size-$2$ candidates
do not satisfy the diagonal closure condition. No size-$2$ subset
is jointly closed.

\emph{Size $3$.} Suppose $S=\{a,b,c\}$ is jointly closed. The
three diagonal entries $B(a,a),B(b,b),B(c,c)$ and the three
$T$-diagonals $T(s,s)$ for $s\in S$ must all lie in $S$. Direct
enumeration of $\binom{10}{3}=120$ candidates against just the
six diagonal constraints leaves only two candidates: $\{0,7,8\}$
and $\{6,7,8\}$. (For $\{0,7,8\}$: $B$-diagonals are $0,8,7$; for
$\{6,7,8\}$: $B$-diagonals are $7,8,7$; the $T$-diagonals are
$0,7,7$ and $7,7,7$ respectively, all in $S$.) Both candidates
fail $B$-closure at the off-diagonal cell $B(7,8)=9$:
$9\notin\{0,7,8\}$ and $9\notin\{6,7,8\}$. No size-$3$ subset is
jointly closed.

\emph{Sizes $4$ through $10$.} At each size, direct enumeration
over the $\binom{10}{k}$ candidates verifies that the listed
$S_{k}$ is the unique jointly-closed subset of its size class
($k\in\{4,5,6,7,8,9\}$); the size-$10$ case is $\Zten$ itself.
The full enumeration of all $1023$ non-empty subsets of $\Zten$,
with the distinguishing failing cell recorded for each
non-jointly-closed subset, is performed by the verification
script \texttt{4core\_verification.py} (\S\ref{sec:verify}).
\end{proof}

\begin{remark}[Structural reading of the chain]\label{rem:sigma-walk}
The chain has a clean structural description in terms of the
permutation $\sigma:\Zten\to\Zten$ defined by
$\sigma=(0)(3)(8)(9)(1\;7\;6\;5\;4\;2)$ in cycle notation. The
operators $\{0,3,8,9\}$ are $\sigma$-fixed; the operators
$\{1,2,4,5,6,7\}$ form a single $6$-cycle. The chain proceeds:
\begin{itemize}\setlength\itemsep{0pt}
\item Step $S_{1}\to S_{4}$: adds $\{7,8,9\}$. Of these, $7$ is
in the $6$-cycle and $\{8,9\}$ are $\sigma$-fixed.
\item Steps $S_{4}\to S_{5}$, $S_{5}\to S_{6}$, $S_{6}\to S_{7}$:
each adds the next element of the $\sigma$-forward orbit of $7$,
namely $\sigma(7)=6$, $\sigma^{2}(7)=5$, $\sigma^{3}(7)=4$.
\item Step $S_{7}\to S_{8}$: adds $3$, the remaining
$\sigma$-fixed operator. The $\sigma$-orbit walk pauses for one
step: $\sigma^{4}(7)=2$ would be the next $\sigma$-orbit element,
but adding $2$ alone forces $3=B(2,2)$ into the closure as well,
producing a size jump of $2$. Adding $3$ alone closes more
cheaply, so $3$ enters first.
\item Steps $S_{8}\to S_{9}$, $S_{9}\to S_{10}$: completes the
$\sigma$-orbit walk, adding $\sigma^{4}(7)=2$ at size $9$ and
$\sigma^{5}(7)=1$ at size $10$.
\end{itemize}
The $\sigma$-fixed operators $\{0,3,8,9\}$ enter at three
positions: $0$ at the bottom (size $1$), $\{8,9\}$ in the
size-$1\to 4$ jump, and $3$ at the size-$7\to 8$ bridge. The
$\sigma$-cycle $\{1,7,6,5,4,2\}$ enters in $\sigma$-forward
order, with the size-$8$ bridge interrupting between
$4=\sigma^{3}(7)$ and $2=\sigma^{4}(7)$. Whether this $\sigma$-walk
structure persists for the analogous pair $(T_{p},B_{p})$ over
$\F_{p}$ for primes $p\in\{2,3,5,7,11,13\}$ is a target of future
work \cite{SandersGish2026FpUniversality}.
\end{remark}

\begin{corollary}[Forbidden sizes]\label{cor:forbidden}
No subset of $\Zten$ of size $2$ or $3$ is jointly closed under
$T$ and $B$.
\end{corollary}

\begin{proof}
Part~(b) of the proof of Theorem~\ref{thm:chain}.
\end{proof}

%% -------- section 4: 4-core --------
\section{The 4-core is jointly closed}\label{sec:4core}

The minimal non-trivial element of the chain in
Theorem~\ref{thm:chain} is the four-element subset
$\Cfour=\{0,7,8,9\}$.

\begin{corollary}[4-core closure]\label{cor:4core}
$\Cfour$ is jointly closed under $T$ and $B$. Explicitly,
$T(\Cfour\times\Cfour)\subseteq\{0,7\}\subset\Cfour$ and
$B(\Cfour\times\Cfour)=\Cfour$.
\end{corollary}

\begin{proof}
Part~(a) of the proof of Theorem~\ref{thm:chain}, displayed
explicitly via the $4\times 4$ restrictions.
\end{proof}

\begin{remark}\label{rem:rank-disparity}
The image of $T\big|_{\Cfour\times\Cfour}$ is the proper subset
$\{0,7\}\subset\Cfour$. The image of $B\big|_{\Cfour\times\Cfour}$
is the full $\Cfour$. The two operations therefore differ
markedly in rank on $\Cfour$. The two restrictions agree on
exactly $4$ of the $16$ cells: $(0,0)\to 0$, $(0,7)\to 7$,
$(7,0)\to 7$, and $(8,8)\to 7$. They disagree on the remaining
$12$ cells. This makes the normalizer identity of
Theorem~\ref{thm:norm} striking, as we now show.
\end{remark}

%% -------- section 5: normalizer --------
\section{The normalizer identity on the 4-core}\label{sec:norm}

For $p\in\R^{10}$ supported on $\Cfour$, write
$(v,h,\beta,r)=(p_{0},p_{7},p_{8},p_{9})$.

\begin{theorem}[Normalizer coincidence]\label{thm:norm}
For every $p\in\R^{10}$ supported on $\Cfour$,
\[
Z_{T}(p)\;=\;Z_{B}(p)\;=\;(v+h+\beta+r)^{2}.
\]
\end{theorem}

\begin{proof}
We compute the four 4-core-supported coordinates of
$\widehat{T}(p)$ and $\widehat{B}(p)$ from
Definition~\ref{def:fuse} and the restricted tables of
Corollary~\ref{cor:4core}.

For $T$: cells of $T\big|_{\Cfour\times\Cfour}$ producing $0$ are
$(0,0),(0,8),(8,0),(0,9),(9,0)$ ($5$ cells). Cells producing $7$
are the remaining $11$ cells. Tabulating fuse coordinates in
$(v,h,\beta,r)$:
\begin{align*}
\widehat{T}(p)_{0}&=v^{2}+2v\beta+2vr,\\
\widehat{T}(p)_{7}&=2vh+h^{2}+2h\beta+2hr+\beta^{2}+2\beta r+r^{2},\\
\widehat{T}(p)_{8}&=0,\qquad\widehat{T}(p)_{9}=0.
\end{align*}
Since $T(\Cfour\times\Cfour)\subseteq\Cfour$
(Corollary~\ref{cor:4core}), $\widehat{T}(p)_{c}=0$ for
$c\notin\Cfour$ when $p$ is $\Cfour$-supported. Summing:
\[
Z_{T}(p)\;=\;v^{2}+h^{2}+\beta^{2}+r^{2}+2(vh+v\beta+vr+h\beta+hr+\beta r)\;=\;(v+h+\beta+r)^{2}.
\]

For $B$: cells of $B\big|_{\Cfour\times\Cfour}$ producing $0$ are
$(0,0),(7,9),(9,7),(9,9)$. Cells producing $7$ are
$(0,7),(7,0),(8,8)$. Cells producing $8$ are
$(0,8),(8,0),(7,7),(8,9),(9,8)$. Cells producing $9$ are
$(0,9),(9,0),(7,8),(8,7)$. Tabulating fuse coordinates:
\begin{align*}
\widehat{B}(p)_{0}&=v^{2}+2hr+r^{2},\\
\widehat{B}(p)_{7}&=2vh+\beta^{2},\\
\widehat{B}(p)_{8}&=2v\beta+h^{2}+2\beta r,\\
\widehat{B}(p)_{9}&=2vr+2h\beta.
\end{align*}
Summing:
\[
Z_{B}(p)\;=\;v^{2}+h^{2}+\beta^{2}+r^{2}+2(vh+v\beta+vr+h\beta+hr+\beta r)\;=\;(v+h+\beta+r)^{2}.
\]
Both equal $(v+h+\beta+r)^{2}$.
\end{proof}

\begin{corollary}[Bilinear simplex normalization]\label{cor:Z=1}
For $p\in\R^{10}$ supported on $\Cfour$ with
$v+h+\beta+r=1$, $Z_{T}(p)=Z_{B}(p)=1$. The convex-combination
operator $\alpha\widehat{T}+(1-\alpha)\widehat{B}$ on
$\Cfour$-supported probability distributions therefore preserves
total mass without explicit re-normalization at any
$\alpha\in[0,1]$.
\end{corollary}

\begin{remark}
The identity $Z_{T}=Z_{B}$ on $\Cfour$ arises by a global
cancellation across the $16$ cells of the restriction, not by
pointwise equality. The two operations agree on only $4$ of the
$16$ cells; the remaining $12$ cells route output values
differently between the four output buckets, in such a way that
the total per-bucket sums coincide. The identity is in fact a
polynomial identity in $(v,h,\beta,r)$, exact and independently
verifiable by symbolic expansion.
\end{remark}

%% -------- section 6: forward directions --------
\section{Forward directions}\label{sec:forward}

The pair $(T,B)$ exhibits substantial additional structure beyond
the chain and normalizer identity established above. We sketch
three concrete extensions, each in preparation as a separate
paper:

\paragraph{Closed-form fixed point and Galois extension
\cite{SandersGish2026FixedPoint}.}
Corollary~\ref{cor:Z=1} shows that the convex-combination
operator $\alpha\widehat{T}+(1-\alpha)\widehat{B}$ acts as a
continuous self-map on the closed convex compact set
$\Delta_{\Cfour}=\{p\in\R^{10}_{\geq 0}:v+h+\beta+r=1\}$ at any
$\alpha\in[0,1]$. At $\alpha=\tfrac12$ the iteration admits a
unique non-degenerate fixed point with explicit closed form in
$\Q(\sqrt{3},y^{*})$, where $y^{*}$ is the positive root of
$y^{2}+(2+\sqrt{3})y-(1+\sqrt{3})=0$ over $\Q(\sqrt{3})$. At this
fixed point, $h^{*}/\beta^{*}=1+\sqrt{3}$ exactly. The companion
ratio $r^{*}/\beta^{*}$ satisfies an irreducible integer quartic
$y^{4}+4y^{3}-y^{2}+2y-2=0$ with Galois group $D_{4}$ over $\Q$,
generating the quartic number field LMFDB \textnumero
$4.2.10224.1$ with field discriminant $-10224=-2^{4}\cdot 3^{2}\cdot 71$,
signature $(2,1)$, and class number $1$. The proof of the closed
form is by elementary algebra and does not require the Brouwer
fixed-point theorem, computer algebra, or numerical iteration.

\paragraph{Mixing-weight observations and an open uniqueness
question \cite{SandersGish2026MixingWeight}.}
At five sample mixing weights
$\alpha\in\{0,\tfrac14,\tfrac12,\tfrac34,1\}$, the runtime
fixed point of $\alpha\widehat{T}+(1-\alpha)\widehat{B}$ on the
full simplex exhibits qualitatively distinct behavior. Only at
$\alpha=\tfrac12$ does the ratio $h^{*}/\beta^{*}$ admit a
small-coefficient quadratic relation; at $\alpha\in\{0,\tfrac14,\tfrac34\}$
no integer-quadratic relation is found at fifty-digit precision
within the coefficient bound $|a|,|b|,|c|\leq 20$. Whether
$\alpha=\tfrac12$ is the unique rational mixing weight in
$(0,1)$ at which the runtime fixed-point ratio is algebraic of
small degree is open.

\paragraph{$\F_p$-universality of the joint-closure structure
\cite{SandersGish2026FpUniversality}.} The construction
underlying $(T,B)$, in which $B$ arises as an instance of the
$\CL_{N}$ family and $T$ as a paired coarse projection, extends
naturally to characteristic $p$ for each prime $p$, producing an
analogous pair $(T_{p},B_{p})$ on $\F_{p}$ (and on
$\F_{p}^{2}\cong\Z/p^{2}\Z$ for small $p$). Whether the chain
rigidity of Theorem~\ref{thm:chain} and the normalizer identity
of Theorem~\ref{thm:norm} persist for $p\in\{2,3,5,7,11,13\}$,
with appropriately analogous shell sizes and absorber/zero
structure, is the question. Preliminary computational evidence
indicates persistence; a structural proof is in preparation.

%% -------- section 7: verification --------
\section{Verification}\label{sec:verify}

The two theorem-level statements of this paper are independently
verified by the script \texttt{4core\_verification.py}, archived
at DOI \texttt{10.5281/zenodo.18852047}. The script implements:

\begin{enumerate}\setlength\itemsep{0pt}
\item \textbf{Joint-closure enumeration} (Theorem~\ref{thm:chain}):
tests joint closure of all $1023$ non-empty subsets of $\Zten$
and verifies the chain structure.
\item \textbf{Normalizer identity} (Theorem~\ref{thm:norm}):
symbolically expands $Z_{T}(p)$ and $Z_{B}(p)$ on $\Cfour$ via
\texttt{sympy} and confirms both equal $(v+h+\beta+r)^{2}$.
\end{enumerate}

The script additionally implements four checks supporting the
forward-direction claims of \S\ref{sec:forward} (closed-form fixed
point at fifty-digit precision, persistence across chain shells,
Galois invariants of the quartic against the LMFDB, and an
$\alpha$-sweep PSLQ at five sample mixing weights). These
additional checks are not theorem-level claims of the present
paper but support the in-preparation companion papers.

The script runs to completion in approximately three seconds on a
modest desktop CPU. The output is deterministic.

%% -------- section 8: scope --------
\section{Scope and limits}\label{sec:scope}

The results of this paper are about the specific tables $T$ and
$B$ of Definition~\ref{def:tables}. We do \emph{not} claim:
\begin{itemize}\setlength\itemsep{0pt}
\item that the chain structure of Theorem~\ref{thm:chain} is
shared by other natural pairs of binary operations on $\Zten$;
\item that the normalizer identity of Theorem~\ref{thm:norm}
arises from a categorical or group-theoretic universal
construction;
\item that the structural reading of Remark~\ref{rem:sigma-walk}
($\sigma$-orbit walk plus $\sigma$-fixed bridge) generalizes to
other pairs $(T,B)$ over $\Zten$ or other moduli; the analogous
question for $(T_{p},B_{p})$ over $\F_{p}$ is the subject of
\cite{SandersGish2026FpUniversality}.
\end{itemize}

We \emph{do} claim that the two theorem-level statements of
\S\ref{sec:chain} and \S\ref{sec:norm} are true of the specific
$T$ and $B$ given.

%% -------- references --------
\begin{thebibliography}{99}

\bibitem{DrapalKepka1985}
A.~Dr\'{a}pal and T.~Kepka.
\newblock Group distances of Latin squares.
\newblock \emph{Commentationes Mathematicae Universitatis Carolinae},
26(2):275--283, 1985.

\bibitem{Kepka1980}
T.~Kepka.
\newblock A note on associative triples of elements in cancellation groupoids.
\newblock \emph{Commentationes Mathematicae Universitatis Carolinae},
21(3):479--487, 1980.

\bibitem{SandersGish2026Sigma}
B.~R.~Sanders and M.~Gish.
\newblock Non-Associativity Decay in Binary Composition Tables over
$\mathbb{Z}/N\mathbb{Z}$.
\newblock Preprint, 2026.

\bibitem{SandersGish2026FixedPoint}
B.~R.~Sanders and M.~Gish.
\newblock A closed-form fixed point and a quartic Galois extension
for a pair of operations on $\mathbb{Z}/10\mathbb{Z}$.
\newblock In preparation, 2026.

\bibitem{SandersGish2026MixingWeight}
B.~R.~Sanders and M.~Gish.
\newblock Mixing-weight observations and an open algebraic uniqueness
question.
\newblock In preparation, 2026.

\bibitem{SandersGish2026FpUniversality}
B.~R.~Sanders and M.~Gish.
\newblock $\mathbb{F}_p$-universality of joint closure for a pair
of binary operations.
\newblock In preparation, 2026.

\end{thebibliography}

\end{document}
